%% 08-conclusion.tex --- included by ../sfc.tex
\section{Conclusion}
\label{sec:conclusion}

\PORTER{} is a deterministic DQC compiler built on a plan-oriented remote
stage: it scores a complete inter-core action on the gate progress it
buys, before any preparation SWAP for it exists, and commits before
preparing. Around that it places a capacity discipline with two relief
tiers, a deterministic initial layout, and a per-circuit choice of
teleport-lookahead horizon.

Three conclusions are supported. First, on the evaluated MQT-Bench suites
\PORTER{} uses substantially fewer ($>40\%$) EPR pairs than published \TeleSABRE{} and completes all twenty-four
instances, including the seven on which \TeleSABRE{} does not converge. Second,
against \dSABRE{} it is competitive at the sizes both target,
$1.02\times$ its pooled EPR at $0.91\times$ its compile time under each
compiler's own protocol, and its advantage grows with the machine,
reaching $0.72\times$ on the $100$--$360$-qubit scale-up series. Third, the ablations attribute these results rather than
leaving them to the system as a whole: capacity relief is what lets a
deterministic-greedy design complete in practice, the initial layout
dominates the difference between the two compilers at the sizes both
target, and the lookahead horizon is what carries \PORTER{} past
\dSABRE{} as the machine grows. The two compilers also
prove complementary rather than ranked: neither wins more than eleven of
the twenty-four instances, and taking the better of the two per instance
reaches $11.6\%$ below \dSABRE{} alone (Section~\ref{sec:discussion}).

One ablation result is negative and delimits the contribution. A richer
gate-teleportation vocabulary does not always pay: at 80 qubits a
vocabulary with gate teleportation removed uses $21.4\%$ fewer pairs than
the full action set, and no more on any circuit it completes. We read
this as the plan scorer undervaluing teleportation's latent return rather
than as a verdict on amortisation. It also points somewhere concrete: a
teleport-only \PORTER{} is both the cheapest vocabulary we measured at
that width and the one \dSABRE{}'s completion proof is written for, which
makes it the natural place to try to earn a guarantee \PORTER{} presently
lacks (Section~\ref{sec:future}).

The evaluation is bounded --- two circuit families at the largest widths,
two non-grid core graphs, horizon rules fitted on the corpus they are
measured on, and a cost model that ignores link contention, latency and
fidelity (\appThreats{}). What follows from those bounds ---
scheduling and fidelity beyond pair count, richer interconnects, a
completion guarantee, and a horizon that scales with the core count ---
is set out in Section~\ref{sec:future}.
